\documentclass{homework}
\course{Math 4182H}
\author{Jim Fowler}
\hwtitle{Honors Analysis II}
\input{preamble}

\begin{document}

\maketitle

\begin{inspiration}
It is good to have an end to journey toward; but it is the journey that matters, in the end.
\byline{Ursula K.\ Le Guin in \textit{The Left Hand of Darkness}}
\end{inspiration}

%%%%%%%%%%%%%%%%%%%%%%%%%%%%%%%%%%%%%%%%%%%%%%%%%%%%%%%%%%%%%%%%
\section{Terminology}

\begin{problem}
  Let $C$ be a smooth curve in $\R^n$ parametrized by $g:[a,b]\to\R^n$ with $g'(t)\neq 0$.
  Define the \textbf{arc length} of $C$.
\end{problem}

\begin{problem}\label{line-integral-definition}%
  Let $F$ be a continuous vector field defined in a neighborhood of a smooth oriented curve $C$ in $\R^n$.
  
  Define the \textbf{line integral} $\displaystyle\int_C F\cdot dx$.
\end{problem}

\begin{problem}\label{conservative-definition}%
  What does it mean for a vector field $F:U\to\R^n$ (with $U\subseteq\R^n$ open) to be \textbf{conservative}?
\end{problem}

\begin{problem}\label{potential-definition}%
  What is a \textbf{potential function} for $F$?
\end{problem}

\begin{problem}
  What is a \textbf{piecewise smooth} curve?  A \textbf{closed} curve?  A \textbf{simple} curve?
\end{problem}


%%%%%%%%%%%%%%%%%%%%%%%%%%%%%%%%%%%%%%%%%%%%%%%%%%%%%%%%%%%%%%%%
\section{Numericals}

\begin{problem}
  One arch of the \textbf{cycloid} is the curve
  $g(t)=(t-\sin t,\; 1-\cos t)$ with $t\in[0,2\pi]$.

  Find the arc length of one arch of the cycloid.
\end{problem}

\begin{problem}
  The \textbf{logarithmic spiral} is the curve
  $g(t) = (e^{-t}\cos t,\; e^{-t}\sin t)$ for $t \in [0, \infty)$.
  Show that this curve has finite total arc length and compute it.
\end{problem}

\begin{problem}
  Define $F : \R^2 \to \R^2$ via $F(x,y)=(x^2 y,\;-xy^2)$. Is $F$ conservative?
\end{problem}

\begin{problem}\label{winding-number}%
  The \textbf{angle form} on $\R^2\setminus\{0\}$ is
  \(
    \omega = \displaystyle\frac{-y\,dx+x\,dy}{x^2+y^2}
    \).

    Integrate $\omega$ over various curves, like
\begin{enumerate}
\item[(a)] the unit circle, oriented counterclockwise,
  \item[(b)] the circle of radius $2$ centered at $(3,0)$, oriented counterclockwise.
  \item[(c)] the unit circle traversed counterclockwise \emph{twice}.
  \end{enumerate}
\end{problem}

\begin{problem}
  Show that on the half-plane $U=\{(x,y):x>0\}$, the form $\omega$ equals $d\!\left(\arctan(y/x)\right)$.
  Conclude that $\omega$ is conservative on $U$, and find a potential.
\end{problem}

%%%%%%%%%%%%%%%%%%%%%%%%%%%%%%%%%%%%%%%%%%%%%%%%%%%%%%%%%%%%%%%%
\section{Exploration}

\begin{problem}\label{cross-partials-necessary}%
  Suppose $U\subseteq\R^n$ is open and $F:U\to\R^n$ is $C^1$ and conservative.

  Give names to the components of $F$ via $F(u) = (F_1(u),F_2(u),\ldots,F_n(u))$.

  Show that
  \(
    \displaystyle\frac{\partial F_i}{\partial x_j}=\frac{\partial F_j}{\partial x_i}\)  on $U$
    for all $i,j$.
\end{problem}

\begin{problem}\label{closed-vs-exact}%
  A $C^1$ vector field $F = (F_1, \ldots, F_n)$ on an open set $U \subseteq \R^n$ is called
  \textbf{closed} if
  $\dfrac{\partial F_i}{\partial x_j} = \dfrac{\partial F_j}{\partial x_i}$ for all $i, j$.
  It is called \textbf{exact} if it is conservative
  (cf.~\ref{conservative-definition}).
  Explain, using \ref{cross-partials-necessary}, why every exact field is closed.
\end{problem}

\begin{problem}\label{ftc-for-line-integrals}%
  Suppose $f:\R^n\to\R$ is $C^1$, and let $C$ be a piecewise smooth curve from $a$ to $b$ in $\R^n$.
  Prove the \textbf{fundamental theorem for line integrals}, namely that
  \[
    \int_C \nabla f\cdot dx = f(b)-f(a).
  \]
\end{problem}

\begin{problem}\label{poincare-lemma}%
  An open set $U\subseteq\R^n$ is \textbf{star-shaped} with respect to the
  origin if for every $x\in U$ and every $t\in[0,1]$, we have $tx\in U$.
  Let $F:U\to\R^n$ be $C^1$ and closed (cf.~\ref{closed-vs-exact}).
  Define $f:U\to\R$ by
  \[
    f(x)=\int_0^1 F(tx)\cdot x\,dt.
  \]
  Prove that $\nabla f = F$.
  This is the \textbf{Poincar\'e lemma}: on star-shaped domains, closed implies exact.
\end{problem}

\begin{problem}\label{integrate-around-boundary}%
  Parametrize the unit circle counterclockwise.
  Compute $\displaystyle\oint x\,dy$ and $\displaystyle \oint y\,dx$.

  Then parametrize the boundary of the rectangle $[0,a]\times[0,b]$ counterclockwise.
  Compute the same two integrals.

  Do the same for the triangle with vertices $(0,0)$, $(1,0)$, $(0,1)$.

  What do you observe? What are these integrals calculating?
\end{problem}

\begin{problem}
  Now consider the ``figure-eight'' curve $\gamma(t)=(\sin t,\;\sin 2t)$ for $t\in[0,2\pi]$.
  
  Compute $\displaystyle\oint x\,dy$.  What does this mean for your observations in \ref{integrate-around-boundary}?
\end{problem}

\begin{problem}
  Prove that arc length is independent of parametrization.
  More precisely, suppose $g:[a,b]\to\R^n$ is $C^1$ with $g'(t)\neq 0$,
  and let $\varphi:[c,d]\to[a,b]$ be $C^1$ with $\varphi'>0$ (an
  orientation-preserving reparametrization).
  Show that $\tilde g=g\circ\varphi$ has the same arc length as $g$.
  What happens if $\varphi'<0$?
\end{problem}

%%%%%%%%%%%%%%%%%%%%%%%%%%%%%%%%%%%%%%%%%%%%%%%%%%%%%%%%%%%%%%%%
\section{Prove or Disprove and Salvage if Possible}

\begin{problem}
  Let $C$ be a piecewise smooth curve in $\R^n$ and $F$ a continuous vector field on $C$.
  Then
  \[
    \left|\int_C F\cdot dx\right| < \int_C |F|\,ds.
  \]
\end{problem}

\begin{problem}
  For any $C^1$ curve $\gamma:[a,b]\to\R^n$, the arc length satisfies \(L(\gamma)\ge|\gamma(b)-\gamma(a)|\),
  with equality if and only if $\gamma$ parametrizes a line segment.
\end{problem}

\end{document}
